% =============================================================================
% Introduction — conservative, evidence-aligned framing
% Drop-in include: % =============================================================================
% Introduction — conservative, evidence-aligned framing
% Drop-in include: % =============================================================================
% Introduction — conservative, evidence-aligned framing
% Drop-in include: % =============================================================================
% Introduction — conservative, evidence-aligned framing
% Drop-in include: \input{sections/intro}
% Target: <= 1.5 pages
% =============================================================================
\section{Introduction}
\label{sec:intro}

Reinforcement learning post-training now underpins much of modern
language-model reasoning and alignment~\citep{ouyang2022training,%
guo2025deepseekr1nature}, but empirical claims in this area are often hard to
compare. PPO~\citep{schulman2017proximal}, GRPO~\citep{shao2024deepseekmath},
and DPO~\citep{rafailov2024direct} are usually evaluated on different model
families, with different framework defaults, different reward functions, and
different notions of ``success.'' What the field needs is not another
leaderboard but a substrate that separates algorithmic conclusions from
implementation, initialization, and task-design effects.

\ours{} is an attempt at such a substrate. It aggregates 79 runs spanning 7 RL
libraries, 5 model families, 32+ checkpoints, and scales from 0.6B to
$\sim$1T parameters, with common reporting conventions and released artifacts.
At the same time, we emphasize a central caveat: the current benchmark is
\emph{not} uniformly mature across tasks. The strongest evidence comes from
short-horizon GSM8K experiments with verifiable binary rewards. Frontier API
runs, tool-use experiments, and code-generation studies are useful case
studies, but many remain single-seed, short-horizon, partially completed, or
custom-evaluated.

The first contribution of the benchmark is therefore methodological rather than
triumphalist: it makes visible how much end-to-end stack choice matters. In our
results, model family, initialization, rollout regime, and framework defaults
can all change the apparent algorithm ranking. Several comparisons that look
like ``PPO vs.\ GRPO'' are better understood as \emph{stack-level comparisons},
because the training backends differ in tokenizer/runtime integration, managed
defaults, and rollout plumbing. This is precisely why a shared measurement
substrate is needed.

Our second contribution is a descriptive diagnostic for sparse-reward GRPO:
Zero-Variance Fraction (ZVF), the fraction of rollout groups with identical
rewards, together with Gradient Utilization (GU = $1-\mathrm{ZVF}$). ZVF is
useful because it directly reveals when within-group contrast disappears.
However, because our main tasks use binary outcome rewards, ZVF is mechanically
coupled to reward sparsity, group size, and baseline accuracy. We therefore use
ZVF/GU as diagnostics of signal degeneracy, not as proof of an independent or
causal failure mode beyond simpler observables such as reward mean, entropy,
advantage variance, or divergence proxies.

Our third contribution is a set of targeted ablations on trainability. In the
current benchmark, instruction-tuned checkpoints are consistently easier to
optimize than comparable base models, and intermediate rollout group sizes
outperform extremes in limited sweeps. These observations support a cautious
claim: whether RL fine-tuning works at all often depends more on
initialization and reward-bearing rollout structure than on the nominal choice
of RL algorithm. They do \emph{not} yet justify a universal optimal group size
or a general ``SFT dominates RL'' law.

Finally, we explicitly separate training reward from held-out generalization.
For Qwen3-8B on GSM8K, a five-seed held-out evaluation yields 83.3\% accuracy,
but the base model under the same greedy protocol already scores 82.0\%; the
+1.3 percentage-point delta is not statistically significant ($p{=}0.26$).
This negative control matters: strong online reward curves do not by themselves
establish generalization gains. We release the benchmark, step-level traces,
and checkpoints so that stronger follow-up work can test exactly these failure
points.

% Target: <= 1.5 pages
% =============================================================================
\section{Introduction}
\label{sec:intro}

Reinforcement learning post-training now underpins much of modern
language-model reasoning and alignment~\citep{ouyang2022training,%
guo2025deepseekr1nature}, but empirical claims in this area are often hard to
compare. PPO~\citep{schulman2017proximal}, GRPO~\citep{shao2024deepseekmath},
and DPO~\citep{rafailov2024direct} are usually evaluated on different model
families, with different framework defaults, different reward functions, and
different notions of ``success.'' What the field needs is not another
leaderboard but a substrate that separates algorithmic conclusions from
implementation, initialization, and task-design effects.

\ours{} is an attempt at such a substrate. It aggregates 79 runs spanning 7 RL
libraries, 5 model families, 32+ checkpoints, and scales from 0.6B to
$\sim$1T parameters, with common reporting conventions and released artifacts.
At the same time, we emphasize a central caveat: the current benchmark is
\emph{not} uniformly mature across tasks. The strongest evidence comes from
short-horizon GSM8K experiments with verifiable binary rewards. Frontier API
runs, tool-use experiments, and code-generation studies are useful case
studies, but many remain single-seed, short-horizon, partially completed, or
custom-evaluated.

The first contribution of the benchmark is therefore methodological rather than
triumphalist: it makes visible how much end-to-end stack choice matters. In our
results, model family, initialization, rollout regime, and framework defaults
can all change the apparent algorithm ranking. Several comparisons that look
like ``PPO vs.\ GRPO'' are better understood as \emph{stack-level comparisons},
because the training backends differ in tokenizer/runtime integration, managed
defaults, and rollout plumbing. This is precisely why a shared measurement
substrate is needed.

Our second contribution is a descriptive diagnostic for sparse-reward GRPO:
Zero-Variance Fraction (ZVF), the fraction of rollout groups with identical
rewards, together with Gradient Utilization (GU = $1-\mathrm{ZVF}$). ZVF is
useful because it directly reveals when within-group contrast disappears.
However, because our main tasks use binary outcome rewards, ZVF is mechanically
coupled to reward sparsity, group size, and baseline accuracy. We therefore use
ZVF/GU as diagnostics of signal degeneracy, not as proof of an independent or
causal failure mode beyond simpler observables such as reward mean, entropy,
advantage variance, or divergence proxies.

Our third contribution is a set of targeted ablations on trainability. In the
current benchmark, instruction-tuned checkpoints are consistently easier to
optimize than comparable base models, and intermediate rollout group sizes
outperform extremes in limited sweeps. These observations support a cautious
claim: whether RL fine-tuning works at all often depends more on
initialization and reward-bearing rollout structure than on the nominal choice
of RL algorithm. They do \emph{not} yet justify a universal optimal group size
or a general ``SFT dominates RL'' law.

Finally, we explicitly separate training reward from held-out generalization.
For Qwen3-8B on GSM8K, a five-seed held-out evaluation yields 83.3\% accuracy,
but the base model under the same greedy protocol already scores 82.0\%; the
+1.3 percentage-point delta is not statistically significant ($p{=}0.26$).
This negative control matters: strong online reward curves do not by themselves
establish generalization gains. We release the benchmark, step-level traces,
and checkpoints so that stronger follow-up work can test exactly these failure
points.

% Target: <= 1.5 pages
% =============================================================================
\section{Introduction}
\label{sec:intro}

Reinforcement learning post-training now underpins much of modern
language-model reasoning and alignment~\citep{ouyang2022training,%
guo2025deepseekr1nature}, but empirical claims in this area are often hard to
compare. PPO~\citep{schulman2017proximal}, GRPO~\citep{shao2024deepseekmath},
and DPO~\citep{rafailov2024direct} are usually evaluated on different model
families, with different framework defaults, different reward functions, and
different notions of ``success.'' What the field needs is not another
leaderboard but a substrate that separates algorithmic conclusions from
implementation, initialization, and task-design effects.

\ours{} is an attempt at such a substrate. It aggregates 79 runs spanning 7 RL
libraries, 5 model families, 32+ checkpoints, and scales from 0.6B to
$\sim$1T parameters, with common reporting conventions and released artifacts.
At the same time, we emphasize a central caveat: the current benchmark is
\emph{not} uniformly mature across tasks. The strongest evidence comes from
short-horizon GSM8K experiments with verifiable binary rewards. Frontier API
runs, tool-use experiments, and code-generation studies are useful case
studies, but many remain single-seed, short-horizon, partially completed, or
custom-evaluated.

The first contribution of the benchmark is therefore methodological rather than
triumphalist: it makes visible how much end-to-end stack choice matters. In our
results, model family, initialization, rollout regime, and framework defaults
can all change the apparent algorithm ranking. Several comparisons that look
like ``PPO vs.\ GRPO'' are better understood as \emph{stack-level comparisons},
because the training backends differ in tokenizer/runtime integration, managed
defaults, and rollout plumbing. This is precisely why a shared measurement
substrate is needed.

Our second contribution is a descriptive diagnostic for sparse-reward GRPO:
Zero-Variance Fraction (ZVF), the fraction of rollout groups with identical
rewards, together with Gradient Utilization (GU = $1-\mathrm{ZVF}$). ZVF is
useful because it directly reveals when within-group contrast disappears.
However, because our main tasks use binary outcome rewards, ZVF is mechanically
coupled to reward sparsity, group size, and baseline accuracy. We therefore use
ZVF/GU as diagnostics of signal degeneracy, not as proof of an independent or
causal failure mode beyond simpler observables such as reward mean, entropy,
advantage variance, or divergence proxies.

Our third contribution is a set of targeted ablations on trainability. In the
current benchmark, instruction-tuned checkpoints are consistently easier to
optimize than comparable base models, and intermediate rollout group sizes
outperform extremes in limited sweeps. These observations support a cautious
claim: whether RL fine-tuning works at all often depends more on
initialization and reward-bearing rollout structure than on the nominal choice
of RL algorithm. They do \emph{not} yet justify a universal optimal group size
or a general ``SFT dominates RL'' law.

Finally, we explicitly separate training reward from held-out generalization.
For Qwen3-8B on GSM8K, a five-seed held-out evaluation yields 83.3\% accuracy,
but the base model under the same greedy protocol already scores 82.0\%; the
+1.3 percentage-point delta is not statistically significant ($p{=}0.26$).
This negative control matters: strong online reward curves do not by themselves
establish generalization gains. We release the benchmark, step-level traces,
and checkpoints so that stronger follow-up work can test exactly these failure
points.

% Target: <= 1.5 pages
% =============================================================================
\section{Introduction}
\label{sec:intro}

Reinforcement learning post-training now underpins much of modern
language-model reasoning and alignment~\citep{ouyang2022training,%
guo2025deepseekr1nature}, but empirical claims in this area are often hard to
compare. PPO~\citep{schulman2017proximal}, GRPO~\citep{shao2024deepseekmath},
and DPO~\citep{rafailov2024direct} are usually evaluated on different model
families, with different framework defaults, different reward functions, and
different notions of ``success.'' What the field needs is not another
leaderboard but a substrate that separates algorithmic conclusions from
implementation, initialization, and task-design effects.

\ours{} is an attempt at such a substrate. It aggregates 79 runs spanning 7 RL
libraries, 5 model families, 32+ checkpoints, and scales from 0.6B to
$\sim$1T parameters, with common reporting conventions and released artifacts.
At the same time, we emphasize a central caveat: the current benchmark is
\emph{not} uniformly mature across tasks. The strongest evidence comes from
short-horizon GSM8K experiments with verifiable binary rewards. Frontier API
runs, tool-use experiments, and code-generation studies are useful case
studies, but many remain single-seed, short-horizon, partially completed, or
custom-evaluated.

The first contribution of the benchmark is therefore methodological rather than
triumphalist: it makes visible how much end-to-end stack choice matters. In our
results, model family, initialization, rollout regime, and framework defaults
can all change the apparent algorithm ranking. Several comparisons that look
like ``PPO vs.\ GRPO'' are better understood as \emph{stack-level comparisons},
because the training backends differ in tokenizer/runtime integration, managed
defaults, and rollout plumbing. This is precisely why a shared measurement
substrate is needed.

Our second contribution is a descriptive diagnostic for sparse-reward GRPO:
Zero-Variance Fraction (ZVF), the fraction of rollout groups with identical
rewards, together with Gradient Utilization (GU = $1-\mathrm{ZVF}$). ZVF is
useful because it directly reveals when within-group contrast disappears.
However, because our main tasks use binary outcome rewards, ZVF is mechanically
coupled to reward sparsity, group size, and baseline accuracy. We therefore use
ZVF/GU as diagnostics of signal degeneracy, not as proof of an independent or
causal failure mode beyond simpler observables such as reward mean, entropy,
advantage variance, or divergence proxies.

Our third contribution is a set of targeted ablations on trainability. In the
current benchmark, instruction-tuned checkpoints are consistently easier to
optimize than comparable base models, and intermediate rollout group sizes
outperform extremes in limited sweeps. These observations support a cautious
claim: whether RL fine-tuning works at all often depends more on
initialization and reward-bearing rollout structure than on the nominal choice
of RL algorithm. They do \emph{not} yet justify a universal optimal group size
or a general ``SFT dominates RL'' law.

Finally, we explicitly separate training reward from held-out generalization.
For Qwen3-8B on GSM8K, a five-seed held-out evaluation yields 83.3\% accuracy,
but the base model under the same greedy protocol already scores 82.0\%; the
+1.3 percentage-point delta is not statistically significant ($p{=}0.26$).
This negative control matters: strong online reward curves do not by themselves
establish generalization gains. We release the benchmark, step-level traces,
and checkpoints so that stronger follow-up work can test exactly these failure
points.
